\documentclass[12pt,a4paper]{article}
\usepackage[margin=2cm]{geometry}
\usepackage{amsmath,amssymb,amsthm,mathtools}
\usepackage{enumitem}
\usepackage{microtype}
\usepackage[hidelinks]{hyperref}
\usepackage{graphicx,pdflscape}
\usepackage{xcolor}

\setlist{nosep}
\setlength{\parindent}{0pt}
\setlength{\parskip}{0.4em}

\newtheorem{theorem}{Theorem}[section]
\newtheorem{lemma}[theorem]{Lemma}
\newtheorem{proposition}[theorem]{Proposition}
\newtheorem{corollary}[theorem]{Corollary}

\theoremstyle{definition}
\newtheorem{definition}[theorem]{Definition}
\newtheorem{example}[theorem]{Example}
\newtheorem{problem}{Exercise}[section]

\theoremstyle{remark}
\newtheorem*{remark}{Remark}

\newenvironment{solution}{\begin{proof}[Solution]}{\end{proof}}

\DeclareMathOperator{\poly}{poly}
\DeclareMathOperator{\negl}{negl}
\DeclareMathOperator{\Adv}{Adv}

\begin{document}
\textbf{To:} Dr. Kristina Host\'{a}kov\'{a},
\href{mailto:kristina.hostakova@inf.ethz.ch}{kristina.hostakova@inf.ethz.ch}\\
\textbf{Subject:} A fully direct proof using a reduction

Dear Dr. Host\'{a}kov\'{a},

Below is a fully direct proof that IND-CPA security implies semantic security,
using a reduction. I also include the alternative final
step by contraposition from Lecture 2, alongside the direct conclusion, to
compare the two presentations.

I previously had the impression that using a reduction necessarily introduced
a contrapositive step in a proof. Writing out the argument below helped me
understand the distinction. My understanding is that a reduction constructs
an algorithm using another algorithm as a subroutine, providing the
interaction that the subroutine expects. Whether we use the resulting
implication directly or by contraposition is a separate choice in presenting
the proof.

Could you let me know whether this understanding is correct?

Best regards,\\
Klevis

\medskip\hrule\medskip
\setcounter{section}{1}
\setcounter{theorem}{13}
\begin{theorem}[Semantic security $\iff$ IND-CPA]
\[
\mathrm{SemanticSecure}(\Pi)
\iff\mathrm{IND\text{-}CPA\text{-}Secure}(\Pi).
\]
\[
\begin{aligned}
\Pi&:\text{ public key encryption scheme}.
\end{aligned}
\]
\end{theorem}

\begin{proof}[By direct proof: IND-CPA $\Longrightarrow$ semantic security]
\leavevmode\par\smallskip\hrule\medskip
\[
\begin{aligned}
\text{Goal:}\quad
&\left[\begin{gathered}
\forall\text{ PPT }D\;
\forall\lambda\in\mathbb{N}:\\
\mathrm{Adv}^{\mathrm{IND\text{-}CPA}}_{\Pi,D}(\lambda)\leq\negl(\lambda)
\end{gathered}\right]\\
&\quad\Longrightarrow\quad
\left[\begin{gathered}
\forall\text{ PPT }A\;\exists\text{ PPT }A'\;
\forall\text{ PPT }\mathsf{Samp}\;\forall\text{ PT }f:\\
\forall\lambda\in\mathbb{N}:\quad
|\epsilon_A(\lambda)-\epsilon_{A'}(\lambda)|\leq\negl(\lambda)
\end{gathered}\right].
\end{aligned}
\]
\[
\begin{aligned}
\text{Idea:}\quad
&\left[\begin{gathered}
\exists\text{ PPT }D\;\exists\text{ polynomial }p\;
\forall\lambda\in\mathbb{N}:\\
|\epsilon_A(\lambda)-\epsilon_{A'}(\lambda)|
\leq p(\lambda)\,\mathrm{Adv}^{\mathrm{IND\text{-}CPA}}_{\Pi,D}(\lambda)
+\negl(\lambda)
\end{gathered}\right]\\
&\quad\overset{\mathrm{IND\text{-}CPA\text{-}Secure}(\Pi)}{\Longrightarrow}\quad
|\epsilon_A-\epsilon_{A'}|\text{ negligible}.
\end{aligned}
\]
\renewcommand{\qedsymbol}{}
\[
\forall\text{ PPT }D:\quad
\underbrace{2\left|P_{\Pi,D,\mathrm{win}}^{\mathrm{IND\text{-}CPA}}(\lambda)-\frac12\right|}_{\mathrm{Adv}_{\Pi,D}^{\mathrm{IND\text{-}CPA}}(\lambda)}
\overset{\mathrm{IND\text{-}CPA\text{-}Secure}(\Pi)}{\leq}\negl(\lambda).
\tag{2}\label{eq:indcpa-negligible-bound}
\]
\[
\text{Idea 2:}\quad\text{Construct the IND-CPA game for Idea 1.}
\]
We want to construct an IND-CPA attacker $D$ so that both
$\epsilon_A$ and $\epsilon_{A'}$ appear in its winning probability.
This would let us relate their difference to $D$'s advantage, as required
by Idea 1. The IND-CPA game already provides a challenge bit $b$ with
two possible values. We use these two branches to include both executions
in the same game. When $b=0$, this runs $A$ on an
encryption of $m$, giving the real execution. When $b=1$, it runs $A$ on
an encryption of $0^{|m|}$, which is exactly how $A'$ is defined below.
Thus the probabilities of correctly computing $f(m,h)$ in the two branches
are $\epsilon_A$ and $\epsilon_{A'}$. We can then use them to calculate
$D$'s winning probability and advantage in this one IND-CPA game,
and thereby bound the semantic-security gap.
\[
\forall\text{ PPT }A:\quad
A'(pk,h,\ell):=A(pk,h,\mathsf{Enc}(pk,0^\ell)).
\]
\[
\begin{aligned}
A'&\text{ depends only on }A\text{ and }\Pi,\\
\forall\text{ PPT }\mathsf{Samp}\;\forall\text{ PT }f:\quad
D&=(D')^A\text{ uses the chosen }\mathsf{Samp}\text{ and }f.
\end{aligned}
\]
\par\pagebreak[0]
\begin{samepage}
\begingroup\small
\[
\begin{array}{c@{\;}c@{\;}|c@{\;}c@{\;}c|}
\multicolumn{2}{c}{\text{IND-CPA challenger}}&\multicolumn{3}{c}{D=(D')^A}\\[3pt]
\cline{3-5}
&&\rule{0pt}{18pt}D'&&A\\[5pt]
\cline{3-5}
&&\multicolumn{3}{|c|}{\rule{0pt}{14pt}D_1}\\[3pt]
\cline{3-5}
(pk,sk)\leftarrow\mathsf{Gen}(1^\lambda)
&\xrightarrow{\ pk\ }&&&\\
&&(m,h)\leftarrow\mathsf{Samp}(pk)&&\\
&&m_0:=m,\quad m_1:=0^{|m|}&&\\
&&m_0\ne m_1&&\\
&\xleftarrow{\ m_0,m_1\ }&&&\\
\cline{3-5}
&&\multicolumn{3}{|c|}{\rule{0pt}{14pt}D_2}\\[3pt]
\cline{3-5}
b\leftarrow_{\$}\{0,1\}&&&&\\
ct\leftarrow\mathsf{Enc}(pk,m_b)
&\xrightarrow{\ ct\ }&&&\\
&&&\xrightarrow{\ pk,h,ct\ }&\\
&&&&x\leftarrow A(pk,h,ct)\\
&&&\xleftarrow{\ x\ }&\\
&&b^*:=\begin{cases}
0,&x=f(m,h),\\
1,&x\ne f(m,h)
\end{cases}&&\\
&\xleftarrow{\ b^*\ }&&&\\
\mathrm{return}\ \mathbf{1}[b^*=b]&&&&\\[3pt]
\cline{3-5}
\end{array}
\]
\endgroup
\end{samepage}
\[
\begin{aligned}
\Pr[D\text{ wins}\mid b=0]
&=\Pr[A(pk,h,\mathsf{Enc}(pk,m))=f(m,h)]\\
&=\epsilon_A(\lambda),\\[4pt]
\Pr[D\text{ wins}\mid b=1]
&=\Pr[A'(pk,h,|m|)\ne f(m,h)]\\
&=1-\Pr[A'(pk,h,|m|)=f(m,h)]\\
&=1-\epsilon_{A'}(\lambda).
\end{aligned}
\]
\[
\begin{aligned}
P_{\Pi,D,\mathrm{win}}^{\mathrm{IND\text{-}CPA}}(\lambda)
&=\Pr[D\text{ wins}\mid b=0]\Pr[b=0]\\
&\quad+\Pr[D\text{ wins}\mid b=1]\Pr[b=1]\\
&=\frac12\epsilon_A(\lambda)+\frac12(1-\epsilon_{A'}(\lambda))\\
&=\frac12+\frac12\bigl(\epsilon_A(\lambda)-\epsilon_{A'}(\lambda)\bigr).
\end{aligned}
\]
\par\medskip\noindent
\begin{minipage}[t][5.6cm][t]{0.48\linewidth}
\vspace{0pt}
{\small\itshape Direct conclusion\par}
\smallskip\hrule\medskip
\centering
\resizebox{\linewidth}{!}{$\displaystyle
\begin{aligned}
|\epsilon_A(\lambda)-\epsilon_{A'}(\lambda)|
&=2\left|P_{\Pi,D,\mathrm{win}}^{\mathrm{IND\text{-}CPA}}(\lambda)-\frac12\right|\\
&\overset{\eqref{eq:indcpa-negligible-bound}}{\leq}\negl(\lambda).
\end{aligned}
$}
\vfill\hrule
\end{minipage}
\hfill
\begin{minipage}[t][5.6cm][t]{0.48\linewidth}
\vspace{0pt}
{\raggedleft\small\itshape By contraposition\par}
\smallskip\hrule\medskip
\centering
\resizebox{\linewidth}{!}{$\displaystyle
\begin{aligned}
|\epsilon_A(\lambda)-\epsilon_{A'}(\lambda)|
&\text{ non-negligible}\\
\Longrightarrow\quad
\Adv_{\Pi,D}^{\mathrm{IND\text{-}CPA}}(\lambda)
&=2\left|P_{\Pi,D,\mathrm{win}}^{\mathrm{IND\text{-}CPA}}(\lambda)-\frac12\right|\\
&=|\epsilon_A(\lambda)-\epsilon_{A'}(\lambda)|
\quad\text{non-negligible}\\
&\overset{\mathrm{IND\text{-}CPA\text{-}Secure}(\Pi)}{\Longrightarrow}\bot.
\end{aligned}
$}
\smallskip
\resizebox{\linewidth}{!}{$\displaystyle
\Longrightarrow\quad
|\epsilon_A(\lambda)-\epsilon_{A'}(\lambda)|\leq\negl(\lambda).
$}
\vfill\hrule
\end{minipage}
\par\medskip
\hrule\medskip
\end{proof}
\end{document}
